\documentclass{beamer}
\usepackage[utf8]{inputenc}

\usetheme{Madrid}
\usecolortheme{default}
\usepackage{amsmath,amssymb,amsfonts,amsthm}
\usepackage{txfonts}
\usepackage{tkz-euclide}
\usepackage{listings}
\usepackage{adjustbox}
\usepackage{array}
\usepackage{tabularx}
\usepackage{gvv}
\usepackage{lmodern}
\usepackage{circuitikz}
\usepackage{tikz}
\usepackage{graphicx}

\setbeamertemplate{page number in head/foot}[totalframenumber]

\usepackage{tcolorbox}
\tcbuselibrary{minted,breakable,xparse,skins}



\definecolor{bg}{gray}{0.95}
\DeclareTCBListing{mintedbox}{O{}m!O{}}{%
  breakable=true,
  listing engine=minted,
  listing only,
  minted language=#2,
  minted style=default,
  minted options={%
    linenos,
    gobble=0,
    breaklines=true,
    breakafter=,,
    fontsize=\small,
    numbersep=8pt,
    #1},
  boxsep=0pt,
  left skip=0pt,
  right skip=0pt,
  left=25pt,
  right=0pt,
  top=3pt,
  bottom=3pt,
  arc=5pt,
  leftrule=0pt,
  rightrule=0pt,
  bottomrule=2pt,
  toprule=2pt,
  colback=bg,
  colframe=orange!70,
  enhanced,
  overlay={%
    \begin{tcbclipinterior}
    \fill[orange!20!white] (frame.south west) rectangle ([xshift=20pt]frame.north west);
    \end{tcbclipinterior}},
  #3,
}
\lstset{
    language=C,
    basicstyle=\ttfamily\small,
    keywordstyle=\color{blue},
    stringstyle=\color{orange},
    commentstyle=\color{green!60!black},
    numbers=left,
    numberstyle=\tiny\color{gray},
    breaklines=true,
    showstringspaces=false,
}
\begin{document}
\title 
{3.2.2}
\date{October 1,2025}
\author 
{V.Sainadh-EE25BTECH11061}
\frame{\titlepage}
\begin{frame}{Question}
Draw a triangle $\triangle ABC$ in which $AB=5\,\text{cm}$, $BC=6\,\text{cm}$ and $\angle ABC=60^\circ$.\\[2pt]
\end{frame}
\begin{frame}{Solution}
$\rightarrow$ Let
\begin{align}
a &= \|\vec{C}-\vec{B}\| = 6\,\text{cm} \tag{1}\\
b &= \|\vec{A}-\vec{C}\| \ \ (\text{to be found}) \tag{2}\\
c &= \|\vec{B}-\vec{A}\| = 5\,\text{cm} \tag{3}
\end{align}

$\rightarrow$ By cosine law in $\triangle ABC$ at $\vec{B}$,
\begin{align}
b^{2} &= a^{2}+c^{2}-2ac\cos B \tag{4}\\
      &= 6^{2}+5^{2}-2\cdot6\cdot5\cdot\cos60^\circ \tag{5}\\
      &= 36+25-60\cdot\frac12 \tag{6}\\
      &= 31 \ \Rightarrow\ b=\sqrt{31}\approx 5.57\,\text{cm} \tag{7}
\end{align}
  \end{frame}
  \begin{frame}{Solution}
  $\rightarrow$ Choose coordinates (placing $B$ at the origin and $BC$ on the $x$-axis):
\begin{align}
\vec{B} &= \myvec{0\\0} \tag{8}\\
\vec{C} &= \myvec{6\\0} \tag{9}\\
\vec{A} &= c\,\myvec{\cos B\\ \sin B}
       = \myvec{\tfrac{5}{2}\\ \tfrac{5\sqrt{3}}{2}}
       \approx \myvec{2.5\\4.33} \tag{10}
\end{align}

$\rightarrow$ Therefore, the vertices of $\triangle ABC$ are:
\begin{align}
\vec{A} &= \myvec{2.5\\4.33}, \qquad
\vec{B} = \myvec{0\\0}, \qquad
\vec{C} = \myvec{6\\0} \tag{11}
\end{align}
\enlargethispage{2\baselineskip}   

\end{frame}

\begin{frame}{Solution}
\begin{figure}[h!]
   \centering
   \includegraphics[width=0.8\columnwidth]{figs/Figure_1.png}
   \caption{Plot of the vector $\vec{r}$}
   \label{Plot_1}
\end{figure}

\end{frame}

 
\begin{frame}[fragile]
\frametitle{C Code }
\begin{lstlisting}
#include <math.h>
int triangle_abc(double AB, double BC, double angleB_deg,
                 double* Ax, double* Ay,
                 double* Bx, double* By,
                 double* Cx, double* Cy,
                 double* b_out)
{
    const double pi = 3.14159265358979323846;
    double B = angleB_deg * (pi/180.0);

    // Place B at origin and C on x-axis
    *Bx = 0.0; *By = 0.0;
    *Cx = BC;  *Cy = 0.0;

    // From B, place A using polar (AB, angle B)
    *Ax = AB * cos(B);
    *Ay = AB * sin(B);


\end{lstlisting}
\end{frame}

\begin{frame}[fragile]
    \frametitle{C Code}
    \begin{lstlisting}
    
    // b = AC
    double dx = *Ax - *Cx, dy = *Ay - *Cy;
    *b_out = sqrt(dx*dx + dy*dy);

    return 0;
}
int triangle_example(double* Ax, double* Ay,
                     double* Bx, double* By,
                     double* Cx, double* Cy,
                     double* b_out)
{
    return triangle_abc(5.0, 6.0, 60.0, Ax, Ay, Bx, By, Cx, Cy, b_out);
}

\end{lstlisting}
\end{frame}


\begin{frame}[fragile]
\frametitle{Python Code}
\begin{lstlisting}
import math
import matplotlib.pyplot as plt

AB = 5.0
BC = 6.0
Bdeg = 60.0
Brad = math.radians(Bdeg)

# Coordinates: place B at origin and C on +x as in your LaTeX
Bx, By = 0.0, 0.0
Cx, Cy = BC, 0.0
Ax, Ay = AB*math.cos(Brad), AB*math.sin(Brad)  # (2.5, 4.330...)

# Optional: AC to confirm = sqrt(31)
AC = math.hypot(Ax-Cx, Ay-Cy)

# Plot
plt.figure(figsize=(5,5))

    \end{lstlisting}
\end{frame}

\begin{frame}[fragile]
    \frametitle{Python Code}
    \begin{lstlisting}
    # triangle edges
plt.plot([Ax, Bx], [Ay, By])  # AB
plt.plot([Bx, Cx], [By, Cy])  # BC
plt.plot([Cx, Ax], [Cy, Ay])  # CA
# points
plt.scatter([Ax, Bx, Cx], [Ay, By, Cy])
plt.text(Ax, Ay, f"  A({Ax:.2f},{Ay:.2f})", ha="left", va="bottom")
plt.text(Bx, By, f"  B({Bx:.2f},{By:.2f})", ha="left", va="top")
plt.text(Cx, Cy, f"  C({Cx:.2f},{Cy:.2f})", ha="left", va="top")

    \end{lstlisting}
\end{frame}

 \begin{frame}[fragile]
    \frametitle{ Python Code}
    \begin{lstlisting}
    plt.title("Triangle ABC (AB=5, BC=6, ∠ABC=60°)")
plt.axis("equal")
plt.grid(True)
plt.xlabel("x")
plt.ylabel("y")
plt.tight_layout()
# Save and/or show
plt.savefig("triangle_abc.png", dpi=200)
plt.show()
print(f"AC ≈ {AC:.4f} (should be sqrt(31) ≈ 5.5678)")
   
\end{lstlisting}
\end{frame}

\begin{frame}[fragile]
\frametitle{callc.py}
    \begin{lstlisting}    
import ctypes
from ctypes import c_double, c_int, byref

lib = ctypes.CDLL("./libtriangle.so")   # adjust path if needed

# void-return (int), args: 11 doubles (3 inputs + 7 outputs)
lib.triangle_abc.argtypes = [c_double, c_double, c_double,
                             ctypes.POINTER(c_double), ctypes.POINTER(c_double),
                             ctypes.POINTER(c_double), ctypes.POINTER(c_double),
                             ctypes.POINTER(c_double), ctypes.POINTER(c_double),
                             ctypes.POINTER(c_double)]
lib.triangle_abc.restype  = c_int

    \end{lstlisting}
\end{frame}
\begin{frame}[fragile]
\frametitle{callc.py}
\begin{lstlisting}
lib.triangle_example.argtypes = [ctypes.POINTER(c_double), ctypes.POINTER(c_double),
                                 ctypes.POINTER(c_double), ctypes.POINTER(c_double),
                                 ctypes.POINTER(c_double), ctypes.POINTER(c_double),
                                 ctypes.POINTER(c_double)]
lib.triangle_example.restype  = c_int

def run_example():
    Ax= c_double(); Ay= c_double()
    Bx= c_double(); By= c_double()
    Cx= c_double(); Cy= c_double()
    b = c_double()

    lib.triangle_example(byref(Ax), byref(Ay), byref(Bx), byref(By),
                         byref(Cx), byref(Cy), byref(b))

\end{lstlisting} 
\end{frame}
\begin{frame}[fragile]
\frametitle{callc.py}
\begin{lstlisting}
     print("Example (AB=5, BC=6, ∠B=60°)")
    print(f"A = ({Ax.value:.2f}, {Ay.value:.2f})")
    print(f"B = ({Bx.value:.2f}, {By.value:.2f})")
    print(f"C = ({Cx.value:.2f}, {Cy.value:.2f})")
    print(f"AC = sqrt(31) ≈ {b.value:.2f}")

def run_general(AB=5.0, BC=6.0, angleB_deg=60.0):
    Ax= c_double(); Ay= c_double()
    Bx= c_double(); By= c_double()
    Cx= c_double(); Cy= c_double()
    b = c_double()

    lib.triangle_abc(AB, BC, angleB_deg,
\end{lstlisting}
\end{frame}

\end{document}